% Mathematical core of the IID implementation supplement.
% Requires amsmath, amssymb, amsthm and environments lemma, theorem, proof.
% Arithmetic lemmas are supplied in a separate fragment.

\section{The original i.i.d. FRLP and its accelerated evaluation}
\label{sec:cert-iid-math}

This section proves that the accelerated backward evaluator constructs a
feasible solution of the fixed-coefficient FRLP in Section~3.3 of the paper.
The mathematical inputs are Lemmas~3.3, A.2, and A.3 of the paper.  The
additional arguments below justify the evaluation of that same finite LP:
they do not introduce stronger reward inequalities or replace any of its
rows.  Arithmetic lower enclosures are established separately below.  In
particular, every comparison used for pruning must use such an enclosure,
rather than a rounded estimate without a proved error bound.

\subsection{Fixed data and literal rows}
\label{sec:cert-iid-rows}

The certificate fixes $N=40$ positive nondecreasing drop rates, $N_t=640$
time intervals, $N_y=1200$ state subdivisions, and $N_q=131093$ rate
intervals.  Write $t_k=k/N_t$ and $\Delta=1/N_t$.  The state knots satisfy
\[
  0=y_0=x_0<y_1=x_1<\cdots<y_{N_y}=x_{N_y}=1,
\]
and the rate endpoints satisfy
\[
  0=r_0<r_1<\cdots<r_{N_q-1}<r_{N_q}=+\infty.
\]
All finite endpoints and drop rates are fixed rational inputs.  On each
state rectangle we use the diagonal from its lower-left to its upper-right
vertex, retaining only the part of the grid in $0\le x\le y\le1$.

Fix a time $t=t_k$, a state $0\le x<y\le1$, and the constant drop rate $b>0$
on $[t,t+\Delta]$.  As in Appendix~A.1, let
\[
 c=\max_{0\le s\le1}\{b(s)+\theta(s)\}-1>0,\qquad
 p=cx+y(1-\theta(t)-b),\qquad u=b(1-y+ty).
\]
For finite $r\ge0$, put
\[
 E(r)=e^{-\Delta r},\qquad
 I_r=\int_0^\Delta e^{-rs}\,ds
     =\begin{cases}(1-E(r))/r,&r>0,\\ \Delta,&r=0.\end{cases}
\]
The successor quantities in Lemma~A.2, viewed as functions of a finite
lower endpoint $r$, are
\begin{align}
 \varphi(r)&=1-y+y\bigl(t+I_r+E(r)(1-t-\Delta)\bigr),\label{eq:cert-iid-phi}\\
 y(r)&=\frac{yE(r)}{\varphi(r)},\label{eq:cert-iid-y}\\
 x(r)&=\min\!\left\{y(r),
   \frac{E(r)\bigl(cx+yb\Delta(r+e^{-r}-1)\bigr)}{c\varphi(r)}\right\}.
 \label{eq:cert-iid-x}
\end{align}
Here $\varphi(r)>0$ and $0\le x(r)\le y(r)\le1$.

For finite $r$, define the parameters of Lemma~A.3 by
\begin{align*}
 C(r)&=\min\{p\Delta,pI_r\},\\
 \alpha(r)&=ybE(r)\Delta^2+(-p)^+I_r\Delta,\\
 \beta(r)&=\Delta(u+ybI_r),\qquad z(r)=\beta(r)e^{-r}.
\end{align*}
Both $\alpha(r)$ and $\beta(r)$ are positive at a non-diagonal state.
The maximum of the two unrestricted bounds in Lemma~A.3 is
\begin{equation}
 \mathcal H(\alpha,\beta)=
 \begin{cases}
 \displaystyle \beta-\frac{\beta^2}{4\alpha},&\beta\le\alpha,\\[4pt]
 \displaystyle \alpha\left(\frac34+\frac12\log\frac\beta\alpha\right),
    &\beta>\alpha.
 \end{cases}
 \label{eq:cert-iid-free}
\end{equation}
Indeed, with $v=\beta/\alpha$, the difference between the quadratic and
logarithmic expressions, divided by $\alpha$, is
\[
 v-\frac{v^2}{4}-\frac34-\frac12\log v.
\]
It vanishes at $v=1$ and has derivative $-(v-1)^2/(2v)\le0$.
This proves the stated choice of branch.  When $z(r)>\alpha(r)/2$, write
\[
 J(r)=z(r)+\frac{\alpha(r)r}{2}
       -\frac{\alpha(r)^2}{24\bigl(z(r)-\alpha(r)/2\bigr)}.
\]
For this fixed time and state, suppress the indices $k,i,j$ in the cost coefficient. For the literal rate interval $[r_{h-1},r_h]$, let $a_h$ be the maximum of
\[
 \min\{p\Delta,0\},\qquad
 C(r_h)+\mathcal H\bigl(\alpha(r_h),\beta(r_h)\bigr),\qquad
 C(r_h)+J(r_h),
 \]
where the last term is included only when its guard holds, and the latter
two terms are omitted when $r_h=+\infty$.  Thus
$a_{N_q}=\min\{p\Delta,0\}$.  These are exactly the constants specified at
the end of Appendix~A.1, before numerical lower enclosures are taken.

Let $v_{k+1,i,j}$ be the already fixed next-layer table, and let
$\widetilde v_{k+1}(y,x)$ denote its ordinary continuous piecewise-affine
interpolant on the stated triangles.  The right-hand side of the literal
row indexed by $h$ is
\begin{equation}
 \mathcal R_h
 =a_h+\varphi(r_{h-1})\,
       \widetilde v_{k+1}\bigl(y(r_{h-1}),x(r_{h-1})\bigr).
 \label{eq:cert-iid-row}
\end{equation}
The current table entry must satisfy $v_{k,i,j}\le\mathcal R_h$ for every
$h$.  The code's zero-based leaf index $h-1$ uses \texttt{A(h-1)} for the
continuation and \texttt{B(h)} for the running cost
(\texttt{verify.cpp}, lines 63--96 and 134).

The evaluator may use a lower bound on $\mathcal R_h$ that is strictly
smaller than the literal expression.  In particular, omitting an
applicable lower bound from the maximum defining $a_h$, imposing a
stronger guard before using $J$, or subtracting a nonnegative allowance
cannot invalidate a candidate that is required to lie below the literal
row.  None of these operations changes the target LP.

\subsection{Lowering the numerical table}
\label{sec:cert-iid-lowering}

The monotonicity of the true functions $V_t$ does not by itself imply
monotonicity of an arbitrary numerical table.  We enforce the property
needed below by lowering the newly computed layer, in the following
order:
\begin{align*}
 &i=1,\ldots,N_y,\quad j=i-1,i-2,\ldots,0:
       &&v_{k,i,j}\leftarrow\min\{v_{k,i,j},v_{k,i,j+1}\};\\
 &i=1,\ldots,N_y,\quad j=0,\ldots,i-1:
       &&v_{k,i,j}\leftarrow\min\{v_{k,i,j},v_{k,i-1,j}\}.
\end{align*}
Diagonal entries remain zero.  These are the two passes in
\texttt{lower\_monotone} (lines 158--160).

\begin{lemma}[Feasibility and monotonicity after lowering]
\label{lem:cert-iid-lowering}
Suppose the next layer is fixed, the current layer is nonpositive with
zero diagonal, and every current literal row is satisfied.  The two
passes preserve those properties and yield
\[
 v_{k,i,j}\le v_{k,i,j+1},\qquad
 v_{k,i,j}\le v_{k,i-1,j}\qquad(0\le j<i\le N_y).
\]
Its interpolant is nonpositive, zero on $x=y$, nondecreasing in $x$, and
nonincreasing in $y$ throughout the triangular state domain.
\end{lemma}
\begin{proof}
Each row constraint for layer $k$ has a current entry only on its left-hand
side; its right-hand side uses the already fixed layer $k+1$.  Decreasing
a current entry therefore preserves every such constraint.  The passes
leave the diagonal unchanged and only decrease nonpositive entries.
They are applied before any earlier layer is constructed, so there are no
already constructed constraints from layer $k-1$ to preserve.

The descending order of $j$ in the first pass makes each row
nondecreasing in $j$.  In the second pass, row $i-1$ has already been
completed.  The new row $i$ is the pointwise minimum of its previous
nondecreasing row and the completed nondecreasing row $i-1$, with the
latter extended by a final zero.  Hence it remains nondecreasing.  Each
column inequality is imposed against the completed preceding row and
cannot be changed by a later update.  This proves both displayed
inequalities.

For the interpolation claim, consider a rectangle with coordinate widths
$d_y=y_{i+1}-y_i>0$ and $d_x=x_{j+1}-x_j>0$.  On its lower triangle with
vertices $(y_i,x_j)$, $(y_{i+1},x_j)$, $(y_{i+1},x_{j+1})$, the two
partial derivatives are
\[
 \partial_x\widetilde v_k
   =\frac{v_{k,i+1,j+1}-v_{k,i+1,j}}{d_x}\ge0,
 \qquad
 \partial_y\widetilde v_k
   =\frac{v_{k,i+1,j}-v_{k,i,j}}{d_y}\le0.
\]
On the upper triangle they are
\[
 \partial_x\widetilde v_k
   =\frac{v_{k,i,j+1}-v_{k,i,j}}{d_x}\ge0,
 \qquad
 \partial_y\widetilde v_k
   =\frac{v_{k,i+1,j+1}-v_{k,i,j+1}}{d_y}\le0.
\]
Only the lower triangle is present when $i=j$.  Continuity across edges
extends these monotonicities along any horizontal or vertical segment in
the domain.  Every interpolated value is a convex combination of
nonpositive entries, and interpolation along the diagonal uses only
zero diagonal entries.
\end{proof}

\subsection{A continuation bound for an entire rate subtree}
\label{sec:cert-iid-continuation}

A tree node consists of consecutive literal intervals.  Write $r_-$ for
its first lower endpoint and $r_+$ for its last upper endpoint; $r_-$ is
finite and $r_+$ may be infinite.  Every lower endpoint occurring in the
node lies in $[r_-,r_+]$.  Define
\[
 \rho(r)=\frac{x}{y}+\frac{b\Delta}{c}(r+e^{-r}-1),
 \qquad x(r)=y(r)\min\{1,\rho(r)\}.
\]

\begin{lemma}[Uniform continuation bound]
\label{lem:cert-iid-continuation}
Suppose $\widetilde v_{k+1}$ has the properties in
Lemma~\ref{lem:cert-iid-lowering}.  If $\rho(r_-)\ge1$, every continuation in
the node is zero.  Otherwise define
\[
 y_+=\begin{cases}y(r_+),&r_+<+\infty,\\0,&r_+=+\infty,\end{cases}
 \qquad x_-=y_+\rho(r_-).
\]
For every finite $r\in[r_-,r_+]$,
\begin{equation}
 \varphi(r)\widetilde v_{k+1}\bigl(y(r),x(r)\bigr)
 \ge
 \varphi(r_-)\widetilde v_{k+1}\bigl(y(r_-),x_-\bigr).
 \label{eq:cert-iid-node-continuation}
\end{equation}
The bound remains valid if $y_+$ is replaced by any certified value in
$[0,y_+]$, in particular by zero.
\end{lemma}
\begin{proof}
The functions $E(r)$ and $I_r$ are nonincreasing.  Since
$1-t-\Delta\ge0$, equation~\eqref{eq:cert-iid-phi} shows that
$\varphi$ is nonincreasing.  Moreover,
\[
 \frac{\varphi(r)}{E(r)}
 =(1-y)e^{\Delta r}
  +y\left(te^{\Delta r}+\int_0^\Delta e^{rs}\,ds+1-t-\Delta\right)
\]
is nondecreasing, so $y(r)=y/(\varphi(r)/E(r))$ is nonincreasing.
Finally $\rho'(r)=b\Delta(1-e^{-r})/c\ge0$.
If $\rho(r_-)\ge1$, it follows that $x(r)=y(r)$, where the interpolant is
zero.

Otherwise $0\le \rho(r_-)<1$.  For every finite $r$ under consideration,
$y(r)\ge y_+$ and $\min\{1,\rho(r)\}\ge \rho(r_-)$; hence
$x(r)\ge x_-$.  Also $y(r)\le y(r_-)$, and
$0\le x_-\le y_+\le y(r_-)$, so the comparison state belongs to the
domain.  Coordinate monotonicity gives
\[
 \widetilde v_{k+1}\bigl(y(r),x(r)\bigr)
 \ge \widetilde v_{k+1}\bigl(y(r_-),x_-\bigr).
\]
Multiplying by $\varphi(r)>0$, and then using
$\varphi(r)\le\varphi(r_-)$ and the nonpositivity of the right-hand
interpolated value, proves~\eqref{eq:cert-iid-node-continuation}.
Replacing $y_+$ by a smaller nonnegative value only decreases $x_-$ and
therefore weakens this lower bound.
\end{proof}

This is the exact-arithmetic argument behind \texttt{rectA}
(lines 98--106).  The implementation uses zero for the far endpoint
unless its arithmetic guard permits an accurately enclosed division.
Its coordinate perturbations and multiplier rounding are justified by
the arithmetic enclosures below; they are not assumptions of exact
monotonicity for rounded expressions.

\subsection{A running-cost bound for an entire rate subtree}
\label{sec:cert-iid-cost}

\begin{lemma}[Monotonicity of the free cost bound]
\label{lem:cert-iid-free}
For fixed $t,x,y,b$, the functions $C(r)$, $\alpha(r)$, and $\beta(r)$
are nonincreasing in finite $r\ge0$.  The function
$\mathcal H(\alpha,\beta)$ is nondecreasing in each positive argument.
Consequently, if $r_+<+\infty$, then every literal cost in the node is
at least
\begin{equation}
 C(r_+)+\mathcal H\bigl(\alpha(r_+),\beta(r_+)\bigr).
 \label{eq:cert-iid-node-free}
\end{equation}
Every literal cost, including the unbounded leaf, is also at least
$\min\{p\Delta,0\}$.
\end{lemma}
\begin{proof}
If $p<0$, then $C(r)=p\Delta$ because $I_r\le\Delta$; if $p\ge0$,
then $C(r)=pI_r$.  Thus $C$ is nonincreasing.  The formulas for $\alpha$
and $\beta$ are nonnegative linear combinations of $E(r)$, $I_r$, and
constants, so the same conclusion holds for them.  For $\beta\le\alpha$,
\[
 \partial_\alpha \mathcal H=\frac{\beta^2}{4\alpha^2}>0,
 \qquad \partial_\beta \mathcal H=1-\frac{\beta}{2\alpha}\ge\frac12.
\]
For $\beta\ge\alpha$,
\[
 \partial_\alpha \mathcal H=\frac14+\frac12\log\frac\beta\alpha>0,
 \qquad \partial_\beta \mathcal H=\frac\alpha{2\beta}>0.
\]
The formulas and their first derivatives agree at $\beta=\alpha$, which
proves coordinate monotonicity on the whole positive quadrant.  Each
finite leaf in the node has upper endpoint at most $r_+$, and its
literal cost includes this free bound evaluated at its own upper
endpoint.  The claimed inequalities follow.
\end{proof}

\begin{lemma}[Uniform shifted cost bound]
\label{lem:cert-iid-shifted}
Suppose $r_+<+\infty$ and $z(r_+)>\alpha(r_-)$.  Then every literal
cost in the node is at least
\begin{equation}
 C(r_+)+z(r_+)+\frac{\alpha(r_+)r_-}{2}
 -\frac{\alpha(r_-)^2}
 {24\bigl(z(r_+)-\alpha(r_-)/2\bigr)}.
 \label{eq:cert-iid-node-shifted}
\end{equation}
\end{lemma}
\begin{proof}
Both factors in $z(r)=\beta(r)e^{-r}$ are nonnegative and
nonincreasing, so $z$ is nonincreasing.  At every finite upper endpoint
$r\in[r_-,r_+]$ of a descendant leaf,
\[
 z(r)\ge z(r_+)>\alpha(r_-)\ge\alpha(r)>\alpha(r)/2.
\]
Thus its literal cost includes the shifted bound $C(r)+J(r)$.
The positive terms satisfy
\[
 C(r)\ge C(r_+),\qquad z(r)\ge z(r_+),\qquad
 \alpha(r)r\ge\alpha(r_+)r_-.
\]
For the subtracted term, its denominator is positive and
\[
 \frac{\alpha(r)^2}{z(r)-\alpha(r)/2}
 \le
 \frac{\alpha(r_-)^2}{z(r_+)-\alpha(r_-)/2}.
\]
Combining these inequalities proves the result.  No monotonicity of the
shifted expression $J(r)$ itself is asserted or required.
\end{proof}

The maximum of the available bounds in the preceding two lemmas is
therefore a valid running-cost lower bound for the node.  For a node
with $r_+=+\infty$ we retain only $\min\{p\Delta,0\}$.  This is the
argument used by \texttt{Bnode} (lines 108--124).  Its stronger numerical
guards, including $r_+\le64$, only suppress optional bounds.  The
implemented free bound uses a conservative logarithm evaluation in
\texttt{log\_lower}; no polynomial replacement of $\mathcal H$ is used or
needed for this certificate.

\subsection{Complete coverage and backward feasibility}
\label{sec:cert-iid-backward}

For each node, add a certified lower enclosure of
\eqref{eq:cert-iid-node-continuation} to a certified lower enclosure of
an applicable running-cost bound.  Call the resulting quantity
$\underline{\mathcal R}_{\rm node}$.  By the preceding lemmas,
\begin{equation}
 \underline{\mathcal R}_{\rm node}\le\mathcal R_h
 \quad\text{for every literal leaf }h\text{ in that node}.
 \label{eq:cert-iid-node-enclosure}
\end{equation}
For an individually evaluated leaf, write
$\underline{\mathcal R}_h\le\mathcal R_h$ for the certified enclosure
returned by its continuation and cost evaluations.

The default run subtracts $2^{-28}$ from each candidate.
It also subtracts a nonnegative quantity of the form
\[
  2^{-40}\left\lceil 2^{24}d_h\right\rceil,
  \qquad 0\le d_h\le y-x\le1,
\]
where the last inequalities are enforced by clamping in
\texttt{allowance} (lines 126--132).  Its detailed formula is only a
candidate-selection policy.  For feasibility, the only fact needed is
that the subtracted quantity is nonnegative.  In particular,
\[
 \underline{\mathcal R}_h-2^{-28}-2^{-40}\lceil 2^{24}d_h\rceil\le\mathcal R_h.
\]
No assertion about a telescoping potential or approximation to an exact
potential decrement is used.

\begin{theorem}[Feasibility of the accelerated backward evaluator]
\label{thm:cert-iid-backward}
Assume the arithmetic routines return the lower enclosures specified
above.  Starting from the zero terminal table, the evaluator's recursive
search and subsequent lowering passes produce a feasible solution of
the original FRLP.
\end{theorem}
\begin{proof}
Fix a state and a previously fixed next layer.  The algorithm maintains a
nonpositive incumbent $w$, initialized by the minimum of zero and
conservative candidate values.  Every update takes a minimum with
$\underline{\mathcal R}_h-2^{-28}-2^{-40}\lceil 2^{24}d_h\rceil$, so $w$ can only decrease.
The initial hint, the unbounded leaf, and the ternary-search proposals
supply candidates only; they do not remove any leaf from the subsequent
search.

A node is pruned only if
$\underline{\mathcal R}_{\rm node}\ge w$.  At that moment,
\eqref{eq:cert-iid-node-enclosure} implies $w\le\mathcal R_h$ for every
leaf in that node.  Later decreases of $w$ preserve all these inequalities.
If a node is not pruned and is a singleton, its leaf is evaluated and the
incumbent is updated; the resulting value is at most that leaf's literal
row.  Otherwise the node is partitioned into two disjoint nonempty
consecutive sets of leaves and both children are searched.

The root contains all $N_q$ leaves, including the last unbounded
interval.  Recursive splitting strictly decreases the number of leaves,
so the search terminates.  Its terminal nodes form a disjoint cover of
the root.  Every terminal node is either pruned with its inequalities
proved, or explicitly evaluated.  Consequently the final incumbent
satisfies $w\le\mathcal R_h$ for every $h$.  This argument also covers an
immediate pruning of the root.  The source's coverage counter, which
adds the number of leaves at a pruned node and one at an evaluated
singleton, checks this complete cover; the counter alone is not a
substitute for~\eqref{eq:cert-iid-node-enclosure}.

Apply this search independently to each non-diagonal state.  Its outputs
are nonpositive and satisfy every literal row against the fixed next
layer; the diagonal is left at zero.  Lemma~\ref{lem:cert-iid-lowering}
then preserves feasibility and supplies the interpolant properties
needed at the next backward step.  The zero terminal table has these
properties.  Backward induction from $k=N_t-1$ to $k=0$ proves the claim.
\end{proof}

The search just analyzed is \texttt{solve}/\texttt{search}
(lines 135--155); the layer update followed by lowering is in
lines 190--201.  In particular, no assumption of unimodality is needed
for the ternary-search proposals, and no value obtained from an ancestor
node is substituted for a literal leaf row.

Finally, the initial table entry is an integer multiple of $2^{-40}$,
and the reported objective is
\[
 \gamma
  =2^{-40}\lfloor2^{40}\theta(0)\rfloor+v_{0,N_y,0}
  \le\theta(0)+v_{0,N_y,0}.
\]
The FRLP implication proved in Section~3.3 of the paper, using the
concavity of the true $V_t$ and Lemma~3.3, therefore gives
\[
 L(b,q)\ge\theta(0)+V_0(1,0)
       \ge\theta(0)+v_{0,N_y,0}
       \ge\gamma
\]
for every admissible $q$.  The numerical interpolant itself need not be
concave.  The all-finite-$n$ consequence is supplied by the paper's
separate transfer theorem; it is not an additional assertion of this
backward evaluator.
